\section{Mệnh đề}

\subsection{Đề 1}
 \Opensolutionfile{ans}[ans/ans-0-B1]
 \caulc
 \begin{ex}%[0D1N1-1]
 	Các kí hiệu nào sau đây dùng để viết đúng mệnh đề \lq\lq $7$ là một số tự nhiên\rq\rq.
 	\choice
 	{$7\subset \mathbb{N}$}
 	{\True $7\in \mathbb{N}$}
 	{$7 < \mathbb{N}$}
 	{$7 \leq \mathbb{N}$}
 	\loigiai{
 		Ký hiệu viết đúng $7\in \mathbb{N}$.
 	}
 \end{ex}
 \begin{ex}%[0D1N1-1]
 	Câu nào sau đây \textbf{không} là mệnh đề?
 	\choice
 	{\True $x > 2$}
 	{$3 < 1$}
 	{$4 - 5 = 1$}
 	{Tam giác đều là tam giác có ba cạnh bằng nhau}
 	\loigiai{
 		Vì $x > 2$ là mệnh đề chứa biến, không phải mệnh đề.
 	}
 \end{ex}	
 \begin{ex}%[0D1N1-1]
 	Với giá trị nào của $x$ thì \lq\lq $x^2-1 = 0, x \in \mathbb{N}$\rq\rq \,là mệnh đề đúng?
 	\choice
 	{\True $x = 1$}
 	{$x = -1$}
 	{$x =\pm 1$}
 	{$x = 0$}
 	\loigiai{
 		Với $x = 1$ thì $1^2-1 = 0$.
 	}
 \end{ex}
 \begin{ex}%[0D1H1-1]
 	Câu nào trong các câu sau \textbf{không} phải là mệnh đề?
 	\choice
 	{$3+2=7$}
 	{$x^2+1 > 0$}
 	{$-2-x^2< 0$}
 	{\True $4+x$}
 	\loigiai{
 		Ta thấy $4+x$ chỉ là một biểu thức, không phải khẳng định.
 	}
 \end{ex}
 \begin{ex}%%[0D1H1-2]
 	Mệnh đề nào sau đây đúng?
 	\choice
 	{\True $\forall x\in \mathbb{R},\,x^2-x+1 > 0$}
 	{$\exists n\in \mathbb{N},\,n < 0$}
 	{$\exists x\in \mathbb{Q},\, x^2=2$}
 	{$\forall x\in \mathbb{Z},\,\dfrac{1}{x}> 0$}
 	\loigiai{
 		Vì $x^2-x+1 = \left(x-\dfrac{1}{2}\right)^2+\dfrac{3}{4} > 0, \forall x \in \mathbb{R}$ .
 	}
 \end{ex}
 
 \begin{ex}%[0D1H1-2]
 	Hỏi trong các mệnh đề sau đây mệnh đề nào là mệnh đề đúng?
 	\choice
 	{\True $\forall x \in \mathbb{R}, x>3 \Rightarrow x^2>9$}
 	{$\forall x \in \mathbb{R}, x>-3 \Rightarrow x^2>9$}
 	{$\forall x \in \mathbb{R}, x^2>9 \Rightarrow x>3$}
 	{$\forall x \in \mathbb{R}, x^2>9 \Rightarrow x>-3$}
 	\loigiai{
 		Ta có \lq\lq$\forall x \in \mathbb{R}, x>3 \Rightarrow x^2>9$\rq\rq\; là mệnh đề đúng.
 	}
 \end{ex}
 
 
 \begin{ex}%[0D1H1-2]
 	Trong các mệnh đề sau, mệnh đề nào đúng?
 	\choice
 	{$6 \sqrt{2}$ là số hữu tỷ}
 	{\True Phương trình $x^2+7x-2=0$ có $2$ nghiệm trái dấu}
 	{$17$ là số chẵn}
 	{Phương trình $x^2+x+7=0$ có nghiệm}
 	\loigiai{
 		Phương trình $x^2+7x-2=0$ có $a\cdot c=-2<0$ nên phương trình có $2$ nghiệm trái dấu.
 	}
 \end{ex}
 
 \begin{ex}%[0D1H1-2]
 	Mệnh đề nào sau là mệnh đề sai?
 	\choice
 	{$\forall n \in \mathbb{N} \colon n \leq 2 n$}
 	{$\exists n \in \mathbb{N} \colon n^2=n$}
 	{\True $\forall x \in \mathbb{R}\colon x^2>0$}
 	{$\exists x \in \mathbb{R}\colon x>x^2$}
 	\loigiai{
 		Mệnh đề $\forall x \in \mathbb{R}\colon x^2>0$ là sai. Vì $\exists 0 \in \mathbb{R}\colon 0^2=0$.
 	}
 \end{ex}
 
 \begin{ex}%[0D1H1-2]
 	Trong các mệnh đề sau tìm mệnh đề đúng?
 	\choice
 	{$\forall x \in \mathbb{R}\colon x^2>0$}
 	{$\forall x \in \mathbb{N}\colon x\vdots 3$}
 	{$\forall x \in \mathbb{R}\colon-x^2<0$}
 	{\True $\exists x \in \mathbb{R}\colon x>x^2$}
 	\loigiai{
 		Mệnh đề $\exists x \in \mathbb{R}\colon x>x^2$ là đúng. Vì $\exists 0{,}5 \in \mathbb{R}\colon 0{,}5<0{,}5^2$.
 	}
 \end{ex}
 
 \begin{ex}%[0D1H1-2]
 	Cho $n$ là số tự nhiên, mệnh đề nào sau đây đúng?
 	\choice
 	{$\forall n,\,n\left(n+1\right)$ là số chính phương}
 	{$\forall n,\,n\left(n+1\right)$ là số lẻ}
 	{$\exists n,\,n\left(n+1\right)\left(n+2\right)$ là số lẻ}
 	{\True $\forall n,\,n\left(n+1\right)\left(n+2\right)$ là số chia hết cho $6$}
 	\loigiai{
 		Mệnh đề $\forall n,\,n\left(n+1\right)\left(n+2\right)$ là tích của $3$ số tự nhiên liên tiếp, trong đó luôn có một số chia hết cho $2$ và một số chia hết cho $3$ nên nó chia hết cho $2\cdot3=6$.
 	}
 \end{ex}
 
 \begin{ex}%[0D1H1-2]
 	Chọn mệnh đề đúng trong các mệnh đề sau:
 	\choice
 	{\True $\forall x\in \mathbb{R}, x+1 > x$}
 	{$\forall x\in \mathbb{R}, \left| x\right|=x$}
 	{$\exists x\in \mathbb{R}, x-3=x^2$}
 	{$\exists x\in \mathbb{R}, x^2< 0$}
 	\loigiai{
 		Mệnh đề $\forall x\in \mathbb{R}, x+1 > x$ luôn đúng.
 	}
 \end{ex}
 
 \begin{ex}%[0D1H1-2]
 	Trong các mệnh đề dưới đây mệnh đề nào đúng?
 	\choice
 	{\True $\forall x \in \mathbb{R}, \mathrm{x}^2+1>0$}
 	{$\forall x \in \mathbb{R}, x^2>x$}
 	{$\exists r \in \mathbb{Q}, r^2=7$}
 	{$\forall n \in \mathbb{N}, n+4$ chia hết cho $4$}
 	\loigiai{
 		Mệnh đề $\forall \mathrm{x} \in \mathbb{R}, \mathrm{x}^2+1>0$ luôn đúng.
 	}
 \end{ex}
 \Closesolutionfile{ans}
 \indapan{6}{ans/ans-0-B1}
 \cauds
 \Opensolutionfile{ans}[ans/ans-0-B1-DS]
 \begin{ex}%[0D1H1-2]
 	Xét tính đúng, sai của các câu sau
 	\choiceTF
 	{\True $P \colon$\lq\lq$3^3$ là số chính phương\rq\rq \,có mệnh đề phủ định là $\overline{P} \colon$\lq\lq$3^{3}$ không là số chính phương\rq\rq}
 	{$Q \colon$\lq\lq Tam giác $A B C$ là tam giác cân\rq\rq\,có mệnh đề phủ định là $\overline{Q}\colon$\lq\lq Tam giác $A B C$ không là tam giác vuông\rq\rq}
 	{\True $R \colon$\lq\lq $2^{2003}-1$ là số nguyên tố\rq\rq\,có mệnh đề phủ định là $\overline{R}\colon$\lq\lq $2^{2003}-1$ không là số nguyên tố\rq\rq}
 	{$H \colon$\lq\lq $\sqrt{2}$ là số vô tỉ\rq\rq \,có mệnh đề phủ định là $\overline{H}\colon\lq\lq \sqrt{2}$ là số hữu tỉ\rq\rq}
 	\loigiai{
 		\begin{itemchoice}
 			\itemch Đúng. Vì $\overline{P}\colon$\lq\lq $3^{3}$ không là số chính phương\rq\rq.
 			\itemch  Sai. Vì $\overline{Q}\colon$\lq\lq Tam giác $A B C$ không là tam giác cân\rq\rq.
 			\itemch Đúng. Vì $\overline{R}\colon$\lq\lq $2^{2003}-1$ không là số nguyên tố\rq\rq.
 			\itemch Sai.
 		\end{itemchoice}
 	}
 \end{ex}
 
 \begin{ex}%[0D1H1-2]
 	Xét tính đúng, sai của các mệnh đề sau
 	\choiceTF
 	{\True $x^{2}-x+1>0$}
 	{\True $24$ chia hết cho $2$ và cho $12$}
 	{$x^{2}+1<0$}
 	{\True $\sqrt{5}$ là số vô tỉ}
 	\loigiai{
 		\begin{itemchoice}
 			\itemch Đúng. Vì $x^{2}-x+1 = (x-\dfrac{1}{2})^2+\dfrac{3}{4}>0$.
 			\itemch Đúng.
 			\itemch Sai. Vì $x^{2}+1<0$ luôn đúng với mọi $x$.
 			\itemch Đúng.
 		\end{itemchoice}
 	}
 \end{ex}
 
 \begin{ex}%[0D1H1-2]
 	Xét tính đúng, sai của các mệnh đề sau
 	\choiceTF
 	{\True $20$ chia hết cho $4$}
 	{\True Tổng hai cạnh trong một tam giác lớn hơn cạnh thứ ba của tam giác đó}
 	{$12$ là một số chính phương}
 	{\True Tích của ba số tự nhiên liên tiếp luôn chia hết cho $3$}
 	\loigiai{
 		\begin{itemchoice}
 			\itemch Đúng. 
 			\itemch Đúng. Vì tổng hai cạnh trong một tam giác lớn hơn cạnh thứ ba của tam giác đó.
 			\itemch Sai. Vì số Số chính phương là số tự nhiên có căn bậc hai là một số tự nhiên.
 			\itemch Đúng.
 		\end{itemchoice}
 	}
 \end{ex}
 
 \begin{ex}%[0D1H1-2]
 	Xét tính đúng, sai của mỗi mệnh đề sau
 	\choiceTF
 	{$\forall x \in \mathbb{R}, x^{2}>0$}
 	{\True $\exists a \in \mathbb{Q}, a>a^{2}$}
 	{\True $\forall n \in \mathbb{Z}, n^{2}+n+2$ chia hết cho $2$}
 	{$\forall n \in \mathbb{N}, n(n+1)(n+2)$ không chia hết cho $3$}
 	\loigiai{
 		\begin{itemchoice}
 			\itemch Mệnh đề sai. Ta chọn $x=0$ thì $x^{2}=0>0$ là sai. 
 			\itemch Mệnh đề đúng. Ta chọn $a=\dfrac{1}{2} \in \mathbb{Q}$ thì $a^{2}=\dfrac{1}{4}$ nên $a>a^{2}$ (đúng).
 			\itemch Mệnh đề đúng. Thật vậy $\forall n \in \mathbb{Z}, n^{2}+n+2=n(n+1)+2$, trong đó $n(n+1)$ là tích của hai số nguyên liên tiếp nên chia hết cho $2$ , vì vậy $n(n+1)+2$ cũng chia hết cho 2.
 			\itemch Mệnh đề sai. Ta cho $n=1$ thì $n(n+1)(n+2)=1 \cdot 2 \cdot 3=6$ chia hết cho $3$.
 		\end{itemchoice}
 	}
 \end{ex}
 
 \Closesolutionfile{ans}
 \indapan{3}{ans/ans-0-B1-DS}
 
 \Opensolutionfile{ans}[ans/ans-0-B1-KQ]
 \caukq
 \begin{ex}%[0D1H1-2]
 	Trong các mệnh đề sau, có bao nhiêu mệnh đề đúng?
 	\begin{enumerate}
 		\item \lq\lq Hình thoi có hai đường chéo vuông góc với nhau\rq\rq.
 		\item \lq\lq$6$ là số nguyên tố\rq\rq.
 		\item \lq\lq Tổng hai cạnh của một tam giác lớn hơn cạnh còn lại\rq\rq.
 		\item \lq\lq Phương trình $x^4-2x^2+2=0$ có nghiệm\rq\rq.
 	\end{enumerate}
 	\shortans{$2$}
 	\loigiai{
 		\begin{enumerate}
 			\item là mệnh đề đúng.
 			\item $6$ không là số nguyên tố. Do đó mệnh đề sai.
 			\item Tổng hai cạnh của một tam giác lớn hơn cạnh còn lại là mệnh đề đúng.
 			\item Phương trình $x^4-2x^2+2=0$ vô nghiệm. Nên là mệnh đề sai.
 		\end{enumerate}
 	}
 \end{ex}
 
 \begin{ex}%[0D1H1-2]
 	Trong các mệnh đề sau, có bao nhiêu mệnh đề đúng?
 	\begin{enumerate}
 		\item \lq\lq$\forall x\in \mathbb{R}, x^2\geq 0$\rq\rq.
 		\item \lq\lq$\text{Tồn tại số tự nhiên đều là số nguyên tố}$\rq\rq.
 		\item \lq\lq$\exists x\in \mathbb{N}, x\  \text{chia hết cho}\ x+1$\rq\rq.
 		\item \lq\lq$\forall x\in \mathbb{N}, n^4-n^2+1\ \text{là hợp số}$\rq\rq.
 	\end{enumerate}
 	\shortans{$3$}
 	\loigiai{
 		\begin{enumerate}
 			\item Mệnh đề đúng.
 			\item Mệnh đề đúng, ví dụ số tự nhiên $2$.
 			\item Mệnh đề đúng, ví dụ $x=0$ thì $0\vdots 1$.
 			\item Mệnh đề sai, ví dụ $n=0$ thì $n^4-n^2+1=1$ không là hợp số.
 		\end{enumerate}
 	}
 \end{ex}
 
 \begin{ex}%[0D1H1-2]
 	Trong các mệnh đề sau, có bao nhiêu mệnh đề sai?
 	\begin{enumerate}
 		\item \lq\lq$\pi <\dfrac{10}{3}$\rq\rq.
 		\item \lq\lq Phương trình $ 3x+7=0$ có nghiệm\rq\rq.
 		\item \lq\lq Có ít nhất một số cộng với chính nó bằng $0$\rq\rq.
 		\item \lq\lq $2022$ là hợp số\rq\rq.
 	\end{enumerate}
 	\shortans{$0$}
 	\loigiai{
 		\begin{enumerate}
 			\item $\pi <\dfrac{10}{3}$.\\
 			Mệnh đề đúng do $\pi\approx 3,14$ và $\dfrac{10}{3}\approx 3,33$nên $\pi <\dfrac{10}{3}$.
 			\item Phương trình $3x+7=0$ có nghiệm.\\
 			Vì phương trình $3x+7=0$ có nghiệm hữu tỉ $ x=-\dfrac{7}{3}$nên mệnh đề là đúng.
 			\item Có ít nhất một số cộng với chính nó bằng 0;Do tồn tại số thực 0 để 0+0=0 nên mệnh đề đúng.
 			\item $2022$ là hợp số.\\
 			Ta có: $2022\rm=\rm1011.2$ nên $2022$ là hợp số hay mệnh đề đã cho là đúng.
 		\end{enumerate}
 	}
 \end{ex}
 
 \begin{ex}%[0D1H1-2]
 	Trong các mệnh đề sau, có bao nhiêu mệnh đề sai?
 	\begin{enumerate}
 		\item \lq\lq Số nguyên tố lớn hơn $2$ là số lẻ\rq\rq.
 		\item \lq\lq Số tự nhiên có chữ số tận cùng là $0$ hoặc $5$ thì chia hết cho $5$\rq\rq.
 		\item \lq\lq Bình phương tất cả các số nguyên đều chia hết cho $2$\rq\rq.
 	\end{enumerate}
 	\shortans{$1$}
 	\loigiai{
 		\begin{enumerate}
 			\item là mệnh đề đúng. Vì mọi số lớn hơn $2$ mà chẵn thì đều chia hết cho $2$, nên không thể là số nguyên tố.
 			\item là mệnh đề đúng.
 			\item là mệnh đề sai. Ví dụ: $3^2=9$ nhưng $9$ không chia hết cho $2$.
 		\end{enumerate}
 	}
 \end{ex}
 \begin{ex}%[0D1H1-2]
 	Trong các mệnh đề sau, có bao nhiêu mệnh đề đúng?
 	\begin{enumerate}
 		\item \lq\lq$\forall x\in \mathbb{R}, x^2\geq 0$\rq\rq.
 		\item  \lq\lq$\text{Tồn tại số tự nhiên đều là số nguyên tố}$\rq\rq.
 		\item \lq\lq$\exists x\in \mathbb{N}, x\  \text{chia hết cho}\ x+1$\rq\rq.
 		\item \lq\lq$\forall x\in \mathbb{N}, n^4-n^2+1\ \text{là hợp số}$\rq\rq.
 	\end{enumerate}
 	\shortans{$3$}
 	\loigiai{
 		\begin{enumerate}
 			\item Mệnh đề đúng.
 			\item  Mệnh đề đúng, ví dụ $2$.
 			\item Mệnh đề đúng, ví dụ $x=0$ thì $0\vdots 1$.
 			\item Mệnh đề sai, ví dụ $n=0$ thì $n^4-n^2+1=1$ không là hợp số.
 		\end{enumerate}
 	}
 \end{ex}
 
 \begin{ex}%[0D1H1-2]
 	Cho mệnh đề chứa biến \lq\lq$P(x)\colon x>x^3$\rq\rq. Trong các mệnh đề sau, có bao nhiêu mệnh đề đúng?
 	\begin{enumerate}
 		\item $P(1)$.
 		\item  $P\left(\dfrac{1}{3}\right)$.
 		\item $\forall x\in \mathbb{N}, P(x)$.
 		\item $\exists x\in \mathbb{N}, \overline{P(x)}$.
 	\end{enumerate}
 	\shortans{$2$}
 	\loigiai{
 		\begin{enumerate}[a)]
 			\item $P(1)$.\\
 			Ta có \lq\lq$P(1)\colon 1>1$\rq\rq \,là mệnh đề sai.
 			\item  $P\left(\dfrac{1}{3}\right)$.\\
 			Ta có \lq\lq$P\left(\dfrac{1}{3}\right)\colon \dfrac{1}{3}>\dfrac{1}{27}$\rq\rq\,là mệnh đề đúng.
 			\item $\forall x\in \mathbb{N}, P(x)$.\\
 			Ta có $\forall x\in \mathbb{N}, P(x)$ \,là mệnh đề sai vì $x=0$ thì \lq\lq$P(0)\colon 0>0$\rq\rq \,là sai.
 			\item $\exists x\in \mathbb{N}, \overline{P(x)}$.\\
 			Ta có $\exists x\in \mathbb{N}, \overline{P(x)}$ \,là mệnh đề đúng, chẳng hạn với $x=2$ thì \lq\lq$P(2)\colon 2\leq 8$\rq\rq \,là đúng.
 		\end{enumerate}
 	}
 \end{ex}
 \Closesolutionfile{ans}
 \indapan{6}{ans/ans-0-B1-KQ}
\newpage 
\subsection{Đề 2}
\Opensolutionfile{ans}[ans/ans-0-B1]
\caulc
\begin{ex}%[0D1H1-4]
	Trong các mệnh đề sau, mệnh đề nào {\bf sai}?
	\choice
	{\True $-\pi <-2 \Leftrightarrow \pi^2 < 4$}
	{$\pi < 4 \Leftrightarrow \pi^2 < 16$}
	{$\sqrt{23} < 5 \Rightarrow 2\cdot \sqrt{23} < 2\cdot 5$}
	{$\sqrt{23} < 5 \Rightarrow-2\cdot \sqrt{23} >-2\cdot 5$}
	\loigiai{
		Mệnh đề kéo theo chỉ sai khi $P$ đúng $Q$ sai.\\
		Vậy mệnh đề ở đáp án $-\pi <-2 \Leftrightarrow \pi^2 < 4$ sai.
	}
\end{ex}
%%%==============EX_2============%%%
\begin{ex}%[0D1H1-5]
	Chọn mệnh đề đúng.
	\choice
	{$\forall n\in {\bf \mathbb{N}}^{*}, n^2-1$ là bội số của $3$}
	{$\exists x\in \mathbb{Q}, x^2=3$}
	{$\forall n\in {\bf \mathbb{N}}, 2^n+1$ là số nguyên tố}
	{\True $\exists n\in {\bf \mathbb{N}}, 2^n \ge n+2$}
	\loigiai{
		Mệnh đề $\exists n\in {\bf \mathbb{N}}, 2^n \ge n+2$ đúng vì tồn tại $2\in {\bf \mathbb{N}}, 2^2 \ge 2+2$.
	}
\end{ex}
%%%==============EX_3============%%%
\begin{ex}%[0D1H1-4]
	Trong các mệnh đề nào sau đây mệnh đề nào {\bf sai}?
	\choice
	{\True Hai tam giác bằng nhau khi và chỉ khi chúng đồng dạng và có một góc bằng nhau}
	{Một tứ giác là hình chữ nhật khi và chỉ khi chúng có $3$ góc vuông}
	{Một tam giác là vuông khi và chỉ khi nó có một góc bằng tổng hai góc còn lại}
	{Một tam giác là đều khi và chỉ khi chúng có hai đường trung tuyến bằng nhau và có một góc bằng $60^{\circ}$}
	\loigiai{ Hai tam giác đồng dạng thì chưa chắc bằng nhau
	}
\end{ex}
%%%==============EX_4============%%%
\begin{ex}%[0D1H1-4]
	Mệnh đề nào sau đây {\bf sai}?
	\choice
	{Tứ giác $ABCD$ là hình chữ nhật $\Rightarrow$ tứ giác $ABCD$ có ba góc vuông}
	{\True Tam giác $ABC$ là tam giác đều $\Leftrightarrow$ $\widehat{A}=60^\circ$}
	{Tam giác $ABC$ cân tại $A$ $\Rightarrow$ $AB=AC$}
	{Tứ giác $ABCD$ nội tiếp đường tròn tâm $O$ $\Rightarrow$ $OA=OB=OC=OD$}
	\loigiai{
	Ý B là hai mệnh đề không tương đương, tam giác $ABC$ có $\widehat{A}=60^\circ$ chưa đủ để nó là tam giác đều.
	}
\end{ex}
%%%==============EX_5============%%%
\begin{ex}%[0D1N1-1]
	Cho mệnh đề chứa biến $P(x)\colon"x+15\le x^2 "$ với $x$ là số thực. Mệnh đề nào sau đây là đúng?
	\choice
	{$P(0)$}
	{$P(3)$}
	{$P(4)$}
	{\True $P(5)$}
	\loigiai{
		$P(5)\colon$ \lq\lq $5+15\le 5^2$\rq\rq.
	}
\end{ex}
%%%==============EX_6============%%%
\begin{ex}%[0D1N2-1]
	Cho biết $x$ là một phần tử của tập hợp $A$, xét các mệnh đề sau:
	\begin{multicols}{4}
		\begin{enumerate}[(I).] 
			\item $x\in A$.
			\item $\{x\}\in A$.
			\item $x\subset A$.
			\item $\{x\}\subset A$.
		\end{enumerate}
	\end{multicols}
	Trong các mệnh đề sau, mệnh đề nào là đúng?
	\choice
	{$I$ và $II$}
	{$I$ và $III$}
	{\True $I$ và $IV$}
	{$II$ và $IV$}
	\loigiai{
		\begin{itemize}
			\item $(II)\colon \{x\}\in A$ sai do giữa hai tập hợp không có quan hệ \lq\lq thuộc\rq\rq.
			\item $(III)\colon x\subset A$ sai do giữa phần tử và tập hợp không có quan hệ \lq\lq con\rq\rq.
		\end{itemize}
	}
\end{ex}
%%%==============EX_7============%%%
\begin{ex}%[0D1H1-1]
	Cho mệnh đề chứa biến $P(n)\colon \lq\lq n^2-1$ chia hết cho $4$\rq\rq\, với $n$ là số nguyên. Xét xem các mệnh đề $P(5)$ và $P(2)$ đúng hay sai?
	\choice
	{$P(5)$ đúng và $P(2)$ đúng}
	{$P(5)$ sai và $P(2)$ sai}
	{\True $P(5)$ đúng và $P(2)$ sai}
	{$P(5)$ sai và $P(2)$ đúng}
	\loigiai{
		$P(5)=24$ đúng do $24\,\vdots\, 4$ còn $P(2)=3$ sai do $3$ không chia hết cho $4$.
	}
\end{ex}
%%%==============EX_8============%%%
\begin{ex}%[0D1H1-4]
	Cho tam giác $ABC$ với $H$ là chân đường cao từ $A$. Mệnh đề nào sau đây {\bf sai}?
	\choice
	{\lq\lq $ABC$ là tam giác vuông ở $A$ $\Leftrightarrow \dfrac{1}{AH^2}=\dfrac{1}{AB^2}+\dfrac{1}{AC^2}$\rq\rq}
	{\lq\lq $ABC$ là tam giác vuông ở $A$ $\Leftrightarrow BA^2=BH.BC$\rq\rq}
	{\lq\lq $ABC$ là tam giác vuông ở $A$ $\Leftrightarrow HA^2=HB.HC$\rq\rq}
	{\True \lq\lq $ABC$ là tam giác vuông ở $A$ $\Leftrightarrow BA^2=BC^2+AC^2$\rq\rq}
	\loigiai{
		Chọn đúng phải là \lq\lq $ABC$ là tam giác vuông ở $A$ $\Leftrightarrow BC^2=AB^2+AC^2$\rq\rq.
	}
\end{ex}
%%%==============EX_9============%%%
\begin{ex}%[0D1H2-1]
	Cho hai số $a=\sqrt{10}+1$, $b=\sqrt{10}-1$. Hãy chọn khẳng định đúng.
	\choice
	{\True $\left(a^2+b^2 \right)\in \mathbb{N}$}
	{\True $\left(a+b\right)\in \mathbb{Q}$}
	{\True $a^2+b^2=20$}
	{\True $a\cdot b=99$}
	\loigiai{
		\begin{itemize}
			\item  Đúng vì $a^2+b^2=22$ là số tự nhiên.
			\item  $2\sqrt{10}$ hiểu nhầm là số hữu tỉ.
			\item  Tính sai $a^2+b^2=11+9=20$.
			\item  Tính sai $a\cdot b=100-1=99$.
		\end{itemize}
	}
\end{ex}
%%%==============EX_10============%%%
\begin{ex}%[0D1C1-6]
	Một tòa nhà có $n$ tầng, các tầng được đánh số từ $1$ đến $n$ theo thứ tự từ dưới lên. Có $4$ thang máy đang ở tầng $1$. Biết rằng mỗi thang máy có thể dừng ở đúng $3$ tầng và $3$ tầng này không là $3$ số nguyên liên tiếp và với hai tầng bất kỳ của tòa nhà luôn có một thang máy dừng được ở cả hang tầng này. Hỏi giá trị lớn nhất của $n$ là bao nhiêu?
	\choice
	{\True $6$}
	{$7$}
	{$8$}
	{$9$}
	\loigiai{
		Giả sử $4$ thang máy đó là $A,B,C,D$.\\
		Do khi bốc hai thang bất kỳ luôn có một thang máy dừng được nên		
		\begin{itemize}
			\item Khi bốc hai tầng $2,3$ có một thang dừng được giả sử đó là thang $A$, nên tầng $4$ không phải thang $A$ dừng.
			\item Khi bốc hai tầng $3,4$ có một thang dừng được giả sử đó là thang $B$, nên tầng $5$ không phải thang $B$ dừng.
			\item Khi bốc hai tầng $4,5$ có một thang dừng được giả sử đó là thang $C$, nên tầng $6$ không phải thang $C$ dừng.
			\item Khi bốc hai tầng $5,6$ có một thang dừng được giả sử đó là thang $D$.
			\item Khi bốc hai tầng $6,7$ có một thang dừng được khi đó không thể là thang $A,B,C$ vì sẽ dừng $4$, thang $D$ không thể ở tầng $7$ do không thể ở ba tầng liên tiếp.
		\end{itemize}
		Vậy khách sạn có tối đa sáu tầng.
	}
\end{ex}
%%%==============EX_11============%%%
\begin{ex}%[0D1V1-5]
	Cho $n$ là số tự nhiên. Mệnh đề nào sau đây đúng?
	\choice
	{\lq\lq $\forall n\in \mathbb{N},n\left(n+1\right)$ là số chính phương\rq\rq}
	{\lq\lq $\forall n\in \mathbb{N},n\left(n+1\right)$ là số lẻ\rq\rq}
	{\lq\lq $\exists n\in \mathbb{N},n\left(n+1\right)\left(n+2\right)$ là số lẻ\rq\rq}
	{\True \lq\lq $\forall n\in \mathbb{N},n\left(n+1\right)\left(n+2\right)$ chia hết cho $6$\rq\rq}
	\loigiai{
		\begin{itemize}
			\item  Với $n=1\Rightarrow n\left(n+1\right)=2$ không phải số chính phương. \\
			$\Rightarrow $ $\forall n\in \mathbb{N},n\left(n+1\right)$ là số chính phương sai.
			\item Với $n=1\Rightarrow n\left(n+1\right)=2$ là số chẵn $\Rightarrow $ $\forall n\in \mathbb{N},n\left(n+1\right)$ là số lẻ sai.
			\item Đặt $P=n\left(n+1\right)\left(n+2\right)$
			\begin{itemize}
				\item Trường hợp 1.  $n$ chẵn $\Rightarrow P$ chẵn.
				\item Trường hợp 2.  $n$ lẻ $\Rightarrow \left(n+1\right)$ chẵn $\Rightarrow P$ chẵn.\\
				Vậy $P$ chẵn $\forall n\in \mathbb{N}$ $\Rightarrow$ $\exists n\in \mathbb{N},n\left(n+1\right)\left(n+2\right)$ là số lẻ sai.			
			\end{itemize}
				\item $P\vdots 6\Leftrightarrow \heva{&P\,\vdots\, 2&(*)\\&P\,\vdots\, 3&(**)}$
			\begin{itemize}
				\item Ở trên ta đã chứng minh $P$ luôn chẵn  $\Rightarrow P\,\vdots\, 2$ nên (*) luôn đúng.
				\item Ta cần chứng minh $P\,\vdots\, 3$.
				\begin{itemize}
					\item Trường hợp 1. $n\,\vdots\, 3$ $\Rightarrow P\,\vdots\, 3$
					\item Trường hợp 2. $n$ chia 3 dư 1 $\Rightarrow \left(n+2\right)\,\vdots\, 3$ $\Rightarrow P\,\vdots\, 3$
					\item Trường hợp 3. $n$ chia 3 dư 2 $\Rightarrow \left(n+1\right)\,\vdots\, 3$ $\Rightarrow P\,\vdots\, 3$.\\
					Vậy $P\,\vdots\, 3$ $\forall n\in \mathbb{N}\Rightarrow P\,\vdots\, 6$.
				\end{itemize}
			\end{itemize}			
		\end{itemize}
	}
\end{ex}
%%%==============EX_12============%%%
\begin{ex}%[0D1V1-6]
	Một nhóm học sinh M, N, P, Q, R xếp thành một hàng dọc trước một quầy nước giải khát. Dưới đây là các thông tin ghi nhận được từ các học sinh trên
	\begin{multicols}{2}
		\begin{itemize}
			\item M, P, R là nam; N, Q là nữ;
			\item M đứng trước Q;
			\item N đứng ở vị trí thứ nhất hoặc thứ hai;
			\item Học sinh đứng sau cùng là nam.
		\end{itemize}
	\end{multicols}
	Hai vị trí nào sau đây phải là hai học sinh khác giới tính?
	\choice
	{Thứ hai và ba}
	{Thứ hai và năm}
	{\True Thứ ba và tư}
	{Thứ ba và năm}
	\loigiai{
		Xét hai trường hợp:
		\begin{itemize}
			\item Trường hợp 1. $N$ đứng ở vị trí thứ nhất.\\
			Do $M$ phải đứng trước $Q$ nên $M$ không thể xếp cuối cùng do đó có các trường hợp xếp:
			\begin{itemize}
				\item $N$ đứng đầu, $M,Q$ đứng tiếp theo và hai bạn sau cùng là $P,R$:
				\begin{center}
					\begin{tabular}{|c|c|c|c|c|}
						\hline
						$N$ & $M$ & $Q$ & $P$ & $R$ \\
						\hline
					\end{tabular} hoặc \begin{tabular}{|c|c|c|c|c|}
						\hline
						$N$ & $M$ & $Q$ & $R$ & $P$ \\
						\hline
					\end{tabular}
				\end{center}			
				Ta loại đáp án \lq\lq Thứ hai và năm\rq\rq.
				\item $N$ đứng đầu, xếp tiếp theo là một bạn nam ($P$ hoặc $R$), $M,Q$, xếp cuối là một bạn nam còn lại
				\begin{center}
					\begin{tabular}{|c|c|c|c|c|}
						\hline
						$N$ & $P$ & $M$ & $Q$ & $R$ \\
						\hline 	
					\end{tabular} hoặc \begin{tabular}{|c|c|c|c|c|}
						\hline
						$N$ & $R$ & $M$ & $Q$ & $P$ \\
						\hline
					\end{tabular}
				\end{center}							
				Ta loại đáp án \lq\lq Thứ ba và năm\rq\rq.
			\end{itemize}
			\item Trường hợp 2. $N$ đứng ở vị trí thứ hai.\\
			Do $M$ phải đứng trước $Q$ nên $M,Q$ phải xếp giữa $N$ và học sinh nam cuối hàng. 
			\begin{center}
				\begin{tabular}{|c|c|c|c|c|}
					\hline
					$R$ & $N$ & $M$ & $Q$ & $P$ \\
					\hline
				\end{tabular} hoặc \begin{tabular}{|c|c|c|c|c|}
					\hline
					$P$ & $N$ & $M$ & $Q$ & $R$ \\
					\hline
				\end{tabular}
			\end{center}			
			Ta loại đáp án \lq\lq Thứ hai và ba\rq\rq.\\
			Vậy hai vị trí thứ ba và tư phải là hai học sinh khác giới tính.
		\end{itemize}
	}
\end{ex}	
\Closesolutionfile{ans}
\indapan{6}{ans/ans-0-B1}
\cauds
\Opensolutionfile{ans}[ans/ans-0-B1-DS]
%%%==============EX_1============%%%
\begin{ex}%[0D1N1-2]
	Xét tính đúng, sai của mỗi mệnh đề sau.
	\choiceTF
	{$15$ không là số nguyên tố}
	{Một tứ giác là hình thoi khi và chỉ khi nó có hai đường chéo vuông góc với nhau}
	{\True$5+19=24$}
	{$6+81=25$}
	\loigiai{
		\begin{itemchoice}
			\itemch là mệnh đề sai.
			\itemch là mệnh đề sai.
			\itemch là mệnh đề đúng.
			\itemch là mệnh đề sai.
		\end{itemchoice}
	}
\end{ex}
%%%==============EX_2============%%%
\begin{ex}%[0D1N1-1]
	Cho mệnh đề chứa biến $P(x)$, xét tính đúng sai của các mệnh đề sau.
	\choiceTF
	{$P(1)$}
	{\True$P\left(\dfrac{1}{3} \right)$}
	{$\forall x\in \mathbb{N},P(x)$}
	{\True$\exists x\in \mathbb{N},P(x)$}
	\loigiai{
		\begin{itemchoice}
			\itemch Ta có $P(1)\colon $ \lq\lq$1> 1^3$\rq\rq\, đây là mệnh đề sai.
			\itemch Ta có $P\left(\dfrac{1}{3} \right)\colon$ \lq\lq$\dfrac{1}{3} > \left(\dfrac{1}{3} \right)^3$\rq\rq\, đây là mệnh đề đúng.
			\itemch Ta có $\forall x\in \mathbb{N},x > x^3$ là mệnh đề sai vì $P(1)$ là mệnh đề sai.
			\itemch Ta có $\exists x\in \mathbb{N},x > x^3$ là mệnh đề đúng vì $P\left(\dfrac{1}{3} \right)$ là mệnh đề đúng.
		\end{itemchoice}
	}
\end{ex}
%%%==============EX_3============%%%
\begin{ex}%[0D1H1-5]
	Xét tính đúng (sai) của các mệnh đề sau
	\choiceTF
	{$\forall x\in \mathbb{R},x^3-x^2+1> 0$}
	{\True$\exists n\in \mathbb{N},n^2+3$ chia hết cho $4$}
	{$P\colon$ \lq\lq $\forall x\in \mathbb{R},\forall y\in \mathbb{R}\colon x+y=1$\rq\rq}
	{\True$Q\colon$\lq\lq $\exists x\in \mathbb{R},\exists y\in \mathbb{R}\colon x+y=2$\rq\rq}
	\loigiai{
		\begin{itemchoice}
			\itemch \lq\lq $\forall x\in \mathbb{R},x^3-x^2+10$\rq\rq  là mệnh đề sai.
			\itemch \lq\lq $\exists n\in \mathbb{N},n^2+3$ chia hết cho $4$\rq\rq là mệnh đề đúng.
			\itemch $P\colon$ \lq\lq $\forall x\in \mathbb{R},\forall y\in \mathbb{R}\colon x+y=1$ \rq\rq là mệnh đề sai.
			Vì với $x=2{,}5;y=1\Rightarrow P(2{,}5;1)$ là mệnh đề sai.
			\itemch $Q\colon$\lq\lq $\exists x\in \mathbb{R},\exists y\in \mathbb{R}\colon x+y=2$\rq\rq là mệnh đề đúng.
			Vì $Q(1{,}5;0{,}5)$ là mệnh đề đúng.
		\end{itemchoice}
	}
\end{ex}
%%%==============EX_4============%%%
\begin{ex}%[0D1H1-2]
	Xét tính đúng (sai) của các mệnh đề sau
	\choiceTF
	{Chiến tranh thế giới lần thứ hai kết thúc năm $1\,946$}
	{Chiến dịch Điện Biên Phủ giành thắng lợi năm $1\,975$}
	{\True Sông Hương chảy qua thành phố Huế}
	{Phố cổ Hội An thuộc tỉnh Quãng Ngãi}
	\loigiai{
		\begin{itemchoice}
			\itemch là mệnh đề sai vì chiến tranh thế giới lần thứ hai kết thúc năm $1\,945$.
			\itemch là mệnh đề sai vì chiến dịch Điện Biên Phủ giành thắng lợi năm $1\,954$.
			\itemch là mệnh đề đúng.
			\itemch là mệnh đề sai vì Phố cổ Hội An thuộc tỉnh Quảng Nam.
		\end{itemchoice}
	}
\end{ex}

\Closesolutionfile{ans}
\indapan{3}{ans/ans-0-B1-DS}

\Opensolutionfile{ans}[ans/ans-0-B1-KQ]
\caukq
%%%==============BT_1==============%%%
\begin{ex}%[0D1C1-4]
	Cho các mệnh đề sau. Hãy cho biết có bao nhiêu mệnh đề đúng?
	\begin{itemize}
		\item $\forall n \in \mathbb{N}, n^2$ chia hết cho $7 \Rightarrow n$ chia hết cho $7$.
		\item $\forall n \in \mathbb{N}, n^2$ chia hết cho $5 \Rightarrow n$ chia hết cho $5$.
		\item  Nếu tam giác $ABC$ không phải là tam giác đều thì tam giác đó có ít nhất một góc nhỏ hơn $60^{\circ}$
		\item $\forall n \in \mathbb{N}, n^2 \vdots 5 \Rightarrow n \vdots 5$.		
	\end{itemize}
	\shortans{$4$}
	\loigiai{
		\begin{itemize}
			\item Ta có thể sử dụng phương pháp chứng minh phản chứng như sau\\
			Giả sử $n$ không chia hết cho $7$, suy ra $n=7m+i$, với $m=0,1,2, \ldots$ và $i=1,2,3,4,5,6$.\\
			Ta có $n^2=49m^2+14i m+i^2$, dễ thấy rằng $i^2$ nhận các giá trị $1$, $4$, $9$, $16$, $25$, $36$ đều không chia hết cho $7$ nên $n^2$ không chia hết cho $7$.
			\item Chứng minh tương tự câu trên.
			\item Giả sử tam giác $ABC$ không phải là tam giác đều và không có góc nào nhỏ hơn $60^{\circ}$, tức là cả ba góc đều lớn hơn hoặc bằng $60^{\circ}$.\\
			Do $\widehat{A} \geq 60^{\circ}$, $\widehat{B} \geq 60^{\circ}$, $\widehat{C} \geq 60^{\circ}$ nên $\widehat{A}+\widehat{B}+\widehat{C} \geq 180^{\circ}$. Mà trong một tam giác, tổng ba góc luôn bằng $180^{\circ}$ hay ta có $\widehat{A}+\widehat{B}+\widehat{C}=180^{\circ}$. Vậy khi đó phải có $\widehat{A}=\widehat{B}=\widehat{C}=60^{\circ}$ hay tam giác $ABC$ đều. Điều này trái với giả thiết.\\
			Vậy nếu tam giác $ABC$ không phải là tam giác đều thì tam giác đó có ít nhất một góc nhỏ hơn $60^{\circ}$.
			\item Để chứng minh mệnh đề đó là đúng, ta dùng phương pháp chứng minh phản chứng. Giả sử tồn tại số tự nhiên $n$ mà $n^2$ chia hết cho  $5$  nhưng $n$ không chia hết cho 5. Khi đó, $n$ có dạng $n=5k \pm 1$ hay $n=5k \pm 2$ với $k \in \mathbb{N}$.\\
			Nếu $n=5k \pm 1$ thì $n^2=(5k \pm 1)^2=25k^2 \pm  $10$  k+1$ không chia hết cho $5$.\\
			Nếu $n=5k \pm 2$ thì $n^2=(5k \pm 2)^2=25k^2 \pm  $20$  k+4$ không chia hết cho $5$.\\
			Điều này trái với giả thiết $n^2$ chia hết cho $5$. Vậy điều giả sử là sai, suy ra \lq\lq Nếu bình phương của một số tự nhiên chia hết cho  $5$  thì số đó chia hết cho $5$\lq\lq\,là mệnh đề đúng.
		\end{itemize}
	}
\end{ex}

%%%==============BT_2==============%%%
\begin{ex}%[0D1H1-1]
	Trong các mệnh đề sau, có bao nhiêu mệnh đề toán học?
	\begin{enumerate}
		\item $A\colon$ \lq\lq Tích hai số thực trái dấu là một số thực âm\rq\rq.
		\item $B\colon$ \lq\lq Mọi số tự nhiên đều là dương\rq\rq.
		\item $C\colon$ \lq\lq Có sự sống ngoài Trái Đất\rq\rq.
		\item $D\colon$ \lq\lq Ngày $1$ tháng $5$ là ngày Quốc tế Lao động\rq\rq.
	\end{enumerate}
	\shortans{$2$}
	\loigiai{
		\begin{enumerate}
			\item Phát biểu \lq\lq Tích hai số thực trái dấu là một số thực âm\rq\rq \, là một mệnh đề toán học.
			\item Phát biểu \lq\lq Mọi số tự nhiên đều là dương\rq\rq \, là một mệnh đề toán học.
			\item Phát biểu \lq\lq Có sự sống ngoài Trái Đất\rq\rq \, không là một mệnh đề toán học (vì không liên quan đến sự kiện Toán học nào).
			\item Phát biểu \lq\lq Ngày 1 tháng 5 là ngày Quốc tế Lao động\rq\rq \, không là một mệnh đề toán học (vì không liên quan đến sự kiện Toán học nào).
		\end{enumerate}
	}
\end{ex}

%%%==============BT_3==============%%%
\begin{ex}%[0D1H1-3]
	Lập mệnh đề phủ định của mỗi mệnh đề sau và cho biết có bao nhiêu mệnh đề phủ định sai.
	\begin{enumerate}
		\item $A\colon$ \lq\lq $\dfrac{5}{1{,}2}$ là một phân số\rq\rq .
		\item $B\colon$ \lq\lq Phương trình $x^2+3x+2=0$ có nghiệm\rq\rq .
		\item $C\colon$ \lq\lq $2^2+2^3=2^{2+3}$\rq\rq .
		\item $D\colon$ \lq\lq Số $2025$ chia hết cho $15$\rq\rq.
	\end{enumerate}
	\shortans{$2$}
	\loigiai{
		\begin{enumerate}
			\item $\overline{A}\colon$ \lq\lq$\dfrac{5}{1,2}$ không là một phân số\rq\rq.\\			
			Đúng vì $\dfrac{5}{1{,}2}$ không là phân số (do $1{,}2$ không là số nguyên)
			\item $\overline{B}\colon$ \lq\lq Phương trình $x^2+3x+2=0$ vô nghiệm\rq\rq.\\			
			Sai vì phương trình $x^2+3x+2=0$ có hai nghiệm là $x=-1$ và $x=-2$.
			\item  $\overline{C}\colon$ \lq\lq $ 2^2+2^3 \ne 2^{2+3}$\rq\rq.\\			
			Đúng vì $2^2+2^3=12\ne 32=2^{2+3}$.
			\item $\overline{D}\colon$ \lq\lq Số 2025 không chia hết cho 15\rq\rq.\\
			Sai vì $2025$ chia hết cho $15$.
		\end{enumerate}		
	}
\end{ex}

%%%==============BT_4==============%%%
\begin{ex}%[0D1H1-4]
	Cho $n$ là số tự nhiên. Xét các mệnh đề:
	\begin{itemize}
		\item $P\colon$ \lq\lq$n$ là một số tự nhiên chia hết cho $16$\rq\rq.
		\item $Q\colon$ \lq\lq$n$ là một số tự nhiên chia hết cho $8$\rq\rq.
	\end{itemize}	
	Cho biết có bao nhiêu mệnh đề đúng trong các mệnh đề \lq\lq $P\Rightarrow Q$\rq\rq; \lq\lq $Q\Rightarrow P$\rq\rq\, và \lq\lq $P\Leftrightarrow Q$\rq\rq
	\shortans{$1$}
	\loigiai{
		\begin{itemize}
			\item Phát biểu mệnh đề $P\Rightarrow Q\colon$\lq\lq Nếu $n$ là một số tự nhiên chia hết cho $16$ thì $n$ là một số tự nhiên chia hết cho $8$\rq\rq.\\		
			Mệnh đề này đúng, vì $n$ chia hết cho 16 thì $n=16\cdot k\left(k\in \mathbb{N}\right)$ thì $n=8\cdot (2k)$ chia hết cho $8$.
			\item Phát biểu mệnh đề $Q\Rightarrow P\colon$\lq\lq Nếu $n$ là một số tự nhiên chia hết cho $8$ thì $n$ là một số tự nhiên chia hết cho $16$\rq\rq \\
			Mệnh đề này sai, chẳng hạn $n=8$ là số tự nhiên chia hết cho 8 nhưng $n$ không chia hết cho $16$.
			\item Do mệnh đề \lq\lq$P\Rightarrow Q\colon$\rq\rq  đúng, mệnh đề \lq\lq$Q\Rightarrow P\colon$\rq\rq  sai nên mệnh đề \lq\lq $P\Leftrightarrow Q$\rq\rq sai.
		\end{itemize}
	}
\end{ex}

%%%==============BT_5==============%%%
\begin{ex}%[0D1H1-2]
	Cho các mệnh đế sau. Hãy cho biết có bao nhiêu mệnh đề sai.
	\begin{itemize}
		\item Phương trình $x^2-3x+8=0$ có nghiệm.
		\item $16$ không là số nguyên tố.
		\item Hai phương trình $x^2-4x+3=0$ và $x^2-\sqrt{x+3}+1=0$ có nghiệm chung.
		\item Buôn Mê Thuột là thành phố của tỉnh Quảng Ngãi.
	\end{itemize}
	\shortans{$2$}
	\loigiai{
		\begin{itemize}
			\item Mệnh đề \textbf{sai}.
			\item Mệnh đề \textbf{đúng} vì  $16$  có thể chia hết cho $1$, $2$, $4$, $8$, $16$.
			\item Mệnh đề \textbf{đúng} vì hai phương trình này có $x=1$ là nghiệm chung (thay $x=1$ vào mỗi phương trình để kiểm chứng).
			\item Mệnh đề \textbf{sai}.
		\end{itemize}
	}
\end{ex}

%%%==============BT_6==============%%%
\begin{ex}%[0D1H1-2]
	Các câu sau đây, câu nào là mệnh đề, có bao nhiêu câu không phải là mệnh đề?
	\begin{itemize}
		\item Trong tam giác tổng ba góc bằng $180^{\circ}$.
		\item $(\sqrt{3}-\sqrt{27})^2$ là số nguyên.
		\item $16$ chia  $3$  dư $1$.
		\item $\sqrt{5}$ là số vô tỉ.
	\end{itemize}
	\shortans{$0$}
	\loigiai{
		\begin{itemize}
			\item \lq\lq Trong tam giác tổng ba góc bằng $180^{\circ}$\rq\rq\,là mệnh đề \textbf{đúng}.
			\item $(\sqrt{3}-\sqrt{27})^2$ là số nguyên là mệnh đề \textbf{đúng}.
			\item Mệnh đề \textbf{đúng}.
			\item Mệnh đề \textbf{đúng}.
		\end{itemize}
	}
\end{ex}
\Closesolutionfile{ans}
\indapan{6}{ans/ans-0-B1-KQ}